% cap_metodologia.tex
\chapter{Materiais e Métodos}
\label{chap:materiais_metodos}

Este capítulo descreve o protocolo experimental utilizado na comparação dos modelos. A metodologia foi definida para preservar a ordem temporal dos dados, evitar vazamento de informação futura e permitir comparação direta entre arquiteturas sob o mesmo contrato de entrada, saída e avaliação.

\section{Fonte de Dados e Escopo}

O estudo utiliza medições horárias da estação Sapo, localizada em Conceição do Mato Dentro, Minas Gerais. Neste trabalho, a sigla CMD é usada apenas como abreviação do município de Conceição do Mato Dentro quando necessário. Os dados fazem parte do monitoramento contínuo de qualidade do ar disponibilizado pela Fundação Estadual do Meio Ambiente de Minas Gerais (FEAM) \cite{feam2026dadosmonitoramento}. A variável-alvo é a concentração de material particulado fino (PM2,5), e o principal poluente auxiliar mantido no contrato final é o material particulado com diâmetro de até 10 micrômetros (PM10). O conjunto original também contém variáveis meteorológicas e gases, mas essas colunas foram removidas do contrato principal porque apresentavam alta esparsidade na estação Sapo e prejudicavam a coerência do conjunto de dados final.

O conjunto consolidado usado nos resultados principais contém 34.991 registros horários, de 2017-01-04 00:30:00 a 2020-12-31 22:30:00. Esse período corresponde ao recorte efetivamente disponível após preparação, criação das variáveis substitutas causais de PM10 e remoção das primeiras janelas necessárias para a formulação supervisionada.

\begin{table}[htb]
    \centering
    \caption{Resumo do contrato experimental principal.}
    \label{tab:contrato_experimental}
    \resizebox{\textwidth}{!}{%
    \begin{tabular}{ll}
        \toprule
        \textbf{Item} & \textbf{Definição} \\
        \midrule
        Estação & Sapo / CMD (Conceição do Mato Dentro, Minas Gerais) \\
        Poluente-alvo & PM2,5 horário \\
        Período do artefato & 2017-01-04 00:30:00 a 2020-12-31 22:30:00 \\
        Número de linhas & 34.991 \\
        Tarefa & 48 horas passadas para prever as próximas 24 horas \\
        Divisão externa & Cronológica por proporção 70/15/15 \\
        Seleção de hiperparâmetros & \textit{Walk-forward} temporal nos 85\% anteriores ao teste \\
        Imputação & Linear no contrato final, com máscaras de valores originalmente ausentes \\
        Variáveis principais & PM2,5, PM10, indicadores de ausência, tempo/Fourier e substitutos causais de PM10 \\
        Avaliação final & Teste externo, sem embaralhamento temporal \\
        \bottomrule
    \end{tabular}%
    }
\end{table}

\section{Análise Exploratória e Escolha da Estação}

A escolha da estação Sapo foi precedida por uma análise exploratória comparativa sobre 38 estações/usinas com PM2,5. Essa etapa combinou critérios de cobertura do alvo, continuidade temporal, presença de PM10 pareado, frequência de eventos altos e equilíbrio das partições cronológicas. O objetivo dessa triagem não foi escolher uma estação apenas por conveniência, mas selecionar um recorte com volume suficiente, variável auxiliar causalmente útil e incidência mínima de eventos de cauda para sustentar a comparação entre modelos.

Nesse diagnóstico, Sapo apareceu como uma das candidatas mais fortes: ficou em segundo lugar na classificação por PM2,5 puro e também em segundo lugar na classificação que inclui contexto de PM10. A estação Rio Doce/Estação Principal obteve pontuação ligeiramente maior, mas teve menor frequência de eventos PM2,5 \(\geq 35\) no diagnóstico preliminar (0,19\%, contra 1,16\% em Sapo) e não pertence ao mesmo foco CMD adotado neste trabalho. Assim, Sapo foi mantida como estação principal por combinar aderência ao escopo, cobertura de PM2,5 de 83,47\% no diagnóstico comparativo, cobertura pareada PM10 de 74,27\% e correlação absoluta PM10--PM2,5 de 0,58.

Essa escolha deve ser lida como um recorte metodológico do estudo, não como afirmação de que Sapo seja a estação mais previsível de toda a rede. O papel da triagem foi selecionar uma estação com dados suficientes, presença de PM10 auxiliar e eventos altos em quantidade mínima para análise de cauda. A generalização para outras estações permanece uma limitação e uma possibilidade de continuidade.

\begin{figure}[htb]
    \centering
    \includegraphics[width=\linewidth]{eda_ranking_estacoes_pm25_pm10.pdf}
    \caption{Classificação preliminar das estações por previsibilidade de PM2,5 com contexto de PM10.}
    \label{fig:eda_ranking_estacoes_pm25_pm10}
\end{figure}

Após a escolha da estação, a análise exploratória foi refeita sobre o conjunto final usado na comparação principal. A Figura \ref{fig:eda_sapo_cobertura_lacunas} mostra que a série não é integralmente contínua: há meses com cobertura elevada, mas também lacunas longas, principalmente na validação. A maior lacuna original de PM2,5 ocorre de 2019-12-11 13:30 a 2020-01-28 18:30, com 48,25 dias, seguida por lacunas de 14,38 dias e 10,83 dias no treino. Por outro lado, ainda existem trechos observados contínuos suficientes para análise de janelas; os maiores blocos contínuos identificados têm cerca de 8,54 dias no treino e 8,13 dias no teste. Esse padrão fundamenta duas decisões metodológicas: manter a série em grade horária regular por interpolação e preservar máscaras de ausência para que valores originalmente imputados não sejam tratados como verdade observada na função de perda nem nas métricas.

\begin{figure}[htb]
    \centering
    \includegraphics[width=\linewidth]{eda_sapo_cobertura_lacunas.pdf}
    \caption{Cobertura mensal e maiores lacunas originais de PM2,5 no artefato Sapo.}
    \label{fig:eda_sapo_cobertura_lacunas}
\end{figure}

A Tabela \ref{tab:eda_sapo_splits} resume a cobertura e a distribuição observada por partição. O teste apresenta a maior cobertura observada de PM2,5 (93,37\%), enquanto a validação concentra a maior proporção de ausências (31,73\%). A distribuição também mostra que os eventos altos são raros em todas as partições: 1,57\% dos pontos observados no treino, 1,45\% na validação e 0,96\% no teste ficam em PM2,5 \(\geq 35\). Essa raridade justifica a análise de cauda no Capítulo \ref{chap:resultados}, pois o MAE global sozinho tende a privilegiar o regime mais frequente e pode esconder suavização em eventos de maior concentração.

\begin{table}[htb]
    \centering
    \caption{Resumo exploratório do artefato Sapo por partição.}
    \label{tab:eda_sapo_splits}
    \resizebox{\textwidth}{!}{%
    \begin{tabular}{lrrrrrr}
        \toprule
        \textbf{Partição} & \textbf{Cobertura PM2,5} & \textbf{Cobertura PM10} & \textbf{Mediana PM2,5} & \textbf{p95} & \textbf{p99} & \textbf{PM2,5 \(\geq 35\)} \\
        \midrule
        Treino & 80,70\% & 81,33\% & 9,00 & 24,00 & 40,00 & 1,57\% \\
        Validação & 68,27\% & 64,50\% & 12,00 & 25,00 & 37,00 & 1,45\% \\
        Teste & 93,37\% & 92,11\% & 12,00 & 24,00 & 34,00 & 0,96\% \\
        \bottomrule
    \end{tabular}%
    }
\end{table}

\begin{figure}[htb]
    \centering
    \includegraphics[width=0.95\linewidth]{eda_sapo_distribuicao_pm25.pdf}
    \caption{Distribuição observada de PM2,5 e frequência de eventos altos por partição.}
    \label{fig:eda_sapo_distribuicao_pm25}
\end{figure}

\subsection{Padrões Temporais de PM2,5 e PM10}

Além da cobertura e da distribuição marginal, o artefato Sapo foi analisado quanto a tendência, persistência e sazonalidade. Essa análise usa apenas pontos originalmente observados nas métricas agregadas, respeitando as máscaras \texttt{PM2,5\_\_miss} e \texttt{PM10\_\_miss}. A Figura \ref{fig:serie_temporal_componentes_sapo} resume visualmente a decomposição descritiva do PM2,5: há uma componente lenta de tendência-ciclo, um perfil horário repetitivo e resíduos com picos abruptos.

\begin{figure}[htb]
    \centering
    \includegraphics[width=\linewidth]{serie_temporal_componentes_sapo.pdf}
    \caption{Série temporal de PM2,5 e decomposição descritiva no artefato Sapo. Elaboração própria com base na decomposição aditiva discutida por \citeonline{hyndman2021forecasting}.}
    \label{fig:serie_temporal_componentes_sapo}
\end{figure}

A Tabela \ref{tab:padroes_temporais_sapo} resume os principais sinais temporais de PM2,5 e PM10. Para PM2,5, a persistência de curto prazo é alta, com autocorrelação de 0,813 em 1 hora, e há repetição diária relevante, com autocorrelação de 0,681 em 24 horas. O sinal semanal existe, mas é mais fraco, com autocorrelação de 0,461 em 168 horas e diferença média entre dias da semana de apenas 0,70. O perfil horário é mais claro: as menores médias de PM2,5 aparecem entre a madrugada e o início da tarde, com mínimo médio próximo de 10,29 às 13h, enquanto o fim da tarde e a noite concentram maiores valores, com máximo médio próximo de 14,77 às 18h.

\begin{table}[htb]
    \centering
    \caption{Padrões temporais observados de PM2,5 e PM10 no artefato Sapo.}
    \label{tab:padroes_temporais_sapo}
    \resizebox{\textwidth}{!}{%
    \begin{tabular}{lrrp{0.44\textwidth}}
        \toprule
        \textbf{Diagnóstico} & \textbf{PM2,5} & \textbf{PM10} & \textbf{Leitura} \\
        \midrule
        Autocorrelação 1h & 0,813 & 0,734 & Forte persistência de curto prazo nos dois poluentes. \\
        Autocorrelação 24h & 0,681 & 0,564 & Repetição diária relevante, justificando defasagens e substitutos causais baseados em 24h. \\
        Autocorrelação 168h & 0,461 & 0,360 & Sinal semanal mais fraco que o diário. \\
        Amplitude média por hora & 4,48 & 11,14 & Ciclo horário claro, mais intenso em PM10. \\
        Amplitude média por dia da semana & 0,70 & 3,68 & Variação semanal pequena diante da variabilidade total. \\
        Amplitude média por mês & 8,52 & 12,95 & Regime sazonal/mensal forte, com elevação em setembro e outubro. \\
        \bottomrule
    \end{tabular}%
    }
\end{table}

O PM10 apresenta comportamento compatível com o PM2,5, mas com amplitude horária maior. Suas menores médias ocorrem por volta de 4h, com valor médio próximo de 20,95, e as maiores médias ocorrem por volta de 20h, com valor médio próximo de 32,09. A correlação de Pearson entre PM2,5 e PM10 observados no mesmo horário foi 0,597, e a correlação de Spearman foi 0,421. Portanto, PM10 é uma variável auxiliar informativa, mas não determina completamente a variação horária de PM2,5.

Também há um padrão mensal importante: PM2,5 e PM10 sobem em setembro e outubro, com média de PM2,5 próxima de 17,97 em outubro e p95 de 44 nesse mês. No entanto, a interpretação como ciclo anual regular deve ser cautelosa, pois o artefato cobre apenas quatro anos, possui lacunas e apresenta mudanças de regime entre anos. A autocorrelação anual aproximada de PM2,5 foi baixa, cerca de 0,038, indicando que a repetição anual é menos estável que a persistência horária e diária. Assim, no protocolo final, os termos temporais e de Fourier são úteis para fornecer informação cíclica ao modelo, mas os picos e resíduos continuam dependendo de fatores não observados diretamente, como meteorologia local, mistura atmosférica, transporte espacial e fontes de emissão.

\clearpage

\section{Formulação da Tarefa}

A tarefa de previsão é multi-horizonte: para cada janela válida, os modelos recebem as 48 horas anteriores e produzem uma previsão para as 24 horas seguintes. Formalmente, para um instante inicial \(t\), a entrada contém observações de \(t-47\) até \(t\), e a saída corresponde aos horizontes \(t+1\) até \(t+24\). Essa formulação foi mantida para todos os modelos da comparação principal e está resumida na Figura \ref{fig:tarefa_48_24_sapo}.

Essa escolha de 48 horas foi reavaliada por uma análise auxiliar de persistência no próprio artefato final de Sapo. A autocorrelação observada de PM2,5 foi alta em uma hora (\(0{,}813\)) e ainda relevante em 24 horas (\(0{,}681\)), confirmando que o ciclo diário carrega informação importante para a previsão. Em 48 horas, a autocorrelação permaneceu moderada (\(0{,}598\)), enquanto o sinal semanal de 168 horas foi mais fraco (\(0{,}461\)). Além disso, linhas de base e modelos lineares diagnósticos com janelas de 1 a 168 horas indicaram que a maior parte do ganho ocorre ao cobrir as primeiras 24 horas, sem melhora consistente ao ampliar a entrada para uma semana completa. Assim, a janela de 48 horas foi mantida como compromisso metodológico: ela preserva dois ciclos diários, permite comparar defasagens de 24 e 48 horas e mantém custo computacional controlado, mas não deve ser interpretada como evidência de que históricos mais longos sejam sempre desnecessários em outros recortes.

\begin{figure}[htb]
    \centering
    \includegraphics[width=0.95\linewidth]{tarefa_48_24_sapo.png}
    \caption{Formulação da tarefa 48 para 24 horas do protocolo Sapo.}
    \label{fig:tarefa_48_24_sapo}
\end{figure}

A divisão do conjunto de dados é cronológica e usa 70\% dos horários para treino final, 15\% para validação de treinamento e 15\% para teste final. A Tabela \ref{tab:split_temporal_sapo} apresenta os limites reais do artefato final, e a Figura \ref{fig:split_70_15_15_sapo} mostra a mesma divisão na linha do tempo. O conjunto de teste contém os horários mais recentes e só é usado para avaliação final.

\begin{table}[htb]
    \centering
    \caption{Divisão cronológica externa do conjunto final de Sapo.}
    \label{tab:split_temporal_sapo}
    \begin{tabular}{lrrl}
        \toprule
        \textbf{Partição} & \textbf{Proporção} & \textbf{Timestamps} & \textbf{Intervalo} \\
        \midrule
        Treino & 70\% & 24.495 & 2017-01-04 00:30 a 2019-10-21 14:30 \\
        Validação & 15\% & 5.248 & 2019-10-21 15:30 a 2020-05-27 06:30 \\
        Teste & 15\% & 5.248 & 2020-05-27 07:30 a 2020-12-31 22:30 \\
        \bottomrule
    \end{tabular}
\end{table}

\begin{figure}[htb]
    \centering
    \includegraphics[width=\linewidth]{split_70_15_15_sapo.png}
    \caption{Linha temporal da divisão cronológica usada no artefato Sapo.}
    \label{fig:split_70_15_15_sapo}
\end{figure}

Neste relatório, portanto, a expressão 70/15/15 se refere à divisão cronológica de avaliação. A seleção de hiperparâmetros não usa o teste final: ela é feita por \textit{walk-forward} dentro dos 85\% anteriores ao teste. Em seguida, os modelos selecionados são treinados no protocolo final usando o trecho de treino e a validação de treinamento para parada antecipada e escolha de ponto de controle. A validação intermediária faz parte do desenvolvimento do modelo; o único conjunto tratado como avaliação independente é o teste final.

\section{Pré-processamento}

O pré-processamento foi desenhado para reduzir vazamento temporal. A imputação é ajustada sobre o trecho de treino e aplicada ao restante do período respeitando a separação temporal. No contrato final, usa-se interpolação linear no bloco de treino; para o bloco posterior ao fim do treino, a transformação não usa âncoras futuras e propaga o último estado imputado disponível. Para cada variável imputada, é mantido um indicador \texttt{\_\_miss}, que informa se o valor original estava ausente.

O alvo PM2,5 também possui uma máscara de observação. Quando essa máscara está disponível, pontos do alvo originalmente imputados são excluídos da função de perda e das métricas observadas, seguindo a distinção entre valor observado e valor ausente ou estimado discutida na literatura de dados faltantes e imputação temporal \cite{rubin1976inference,cao2018brits,du2023saits}. Assim, o modelo pode receber uma série regularizada para formar janelas, mas não é recompensado por prever como ``verdade'' um valor que foi produzido pelo próprio pré-processamento.

A normalização é calculada com estatísticas do conjunto de treino. No treinamento dos modelos neurais, médias e desvios são computados por estação e aplicados às janelas de treino, validação e teste. Essa decisão evita que informação agregada dos períodos de validação e teste entre no ajuste de escala do modelo.

\section{Engenharia de Atributos}

O conjunto final de atributos foi mantido enxuto. O codificador recebe PM2,5, PM10, indicadores de ausência e atributos temporais. O decodificador dos modelos que aceitam entradas futuras recebe apenas variáveis conhecidas causalmente no momento da previsão: termos temporais/Fourier e substitutos derivados de PM10. A Tabela \ref{tab:variaveis_contrato} resume os grupos de variáveis.

\begin{table}[htb]
    \centering
    \caption{Variáveis do protocolo final.}
    \label{tab:variaveis_contrato}
    \resizebox{\textwidth}{!}{%
    \begin{tabular}{lll}
        \toprule
        \textbf{Grupo} & \textbf{Variáveis} & \textbf{Uso} \\
        \midrule
        Alvo e histórico & \texttt{PM2,5} & Entrada histórica e alvo da previsão \\
        Poluente auxiliar & \texttt{PM10} & Entrada histórica no codificador \\
        Ausência & \texttt{PM2,5\_\_miss}, \texttt{PM10\_\_miss} & Indicar valores originalmente ausentes \\
        Tempo & \texttt{hour}, \texttt{day\_of\_week}, \texttt{month} & Sazonalidade discreta \\
        Fourier & termos seno/cosseno de 6h, 12h, 24h, 168h e 8760h & Sazonalidade cíclica \\
        Substituto PM10 & \texttt{PM10\_proxy\_exante} & PM10 do mesmo horário no dia anterior, isto é, PM10(t-24) \\
        Tendência PM10 & \texttt{PM10\_trend\_exante} & Diferença causal PM10(t-24) - PM10(t-27) \\
        \bottomrule
    \end{tabular}%
    }
\end{table}

O uso de PM10 no decodificador não assume conhecimento futuro de PM10 observado. O substituto causal usa o valor defasado em 24 horas, e a tendência causal usa a diferença entre valores também defasados. Portanto, ambas podem ser calculadas no instante de emissão da previsão. Variáveis meteorológicas e gases muito esparsos foram excluídos do contrato principal; sua reintrodução deve ser tratada apenas como nova ablação, não como parte da comparação final.

\section{Modelos Avaliados}

Foram comparadas cinco variantes neurais e um modelo de referência baseado em árvores. A Tabela \ref{tab:modelos_metodologia} resume o papel de cada modelo no desenho experimental.

\begin{table}[htb]
    \centering
    \caption{Modelos avaliados na comparação principal.}
    \label{tab:modelos_metodologia}
    \resizebox{\textwidth}{!}{%
    \begin{tabular}{lll}
        \toprule
        \textbf{Modelo} & \textbf{Descrição} & \textbf{Papel na comparação} \\
        \midrule
        LSTM recursiva & Prevê passo a passo, realimentando a trajetória prevista & Linha de base sequencial recursiva \\
        LSTM direta & Produz os 24 horizontes de uma vez & Modelo sequencial direto \\
        Seq2Seq básico & Codificador-decodificador recorrente sem atenção explícita & Linha de base codificador-decodificador \\
        Seq2Seq atenção canônica & Codificador-decodificador com mecanismo de atenção e perda Huber & Variante de referência para análise de atenção \\
        Seq2Seq atenção novo & Codificador-decodificador com atenção, perda ponderada, reinício do decodificador e agrupamento por defasagens & Variante final de Seq2Seq \\
        XGBoost com saída múltipla & Árvores com saída multi-horizonte & Referência baseada em árvores sobre atributos de janelas \\
        \bottomrule
    \end{tabular}%
    }
\end{table}

A escolha dos modelos também delimita o escopo do estudo. Transformers temporais, como o \textit{Temporal Fusion Transformer} (TFT), Informer, Autoformer e PatchTST, são alternativas modernas para previsão multi-horizonte \cite{lim2021tft,zhou2021informer,wu2021autoformer,nie2023patchtst}. Eles não foram incluídos na comparação final porque uma comparação justa exigiria nova família de arquitetura, espaço próprio de hiperparâmetros, orçamento adicional de otimização de hiperparâmetros (HPO), avaliação com múltiplas sementes e decisões específicas de regularização para o tamanho e a qualidade do conjunto de dados Sapo. Além disso, a literatura não autoriza assumir que Transformers vencem automaticamente em séries temporais; \citeonline{zeng2023transformers} mostram que linhas de base simples podem superar Transformers em alguns cenários de previsão. Assim, o recorte deste TCC é comparar famílias recorrentes e modelos baseados em atributos derivados das janelas temporais sob um contrato causal rastreável, deixando Transformers como extensão natural para trabalhos futuros.

A comparação com XGBoost exige uma ressalva metodológica. O modelo de referência baseado em árvores usa o mesmo contrato causal de entrada e é avaliado no mesmo bloco final de teste, mas seu treinamento com saída múltipla exige janelas com horizonte completo observado. Nos modelos neurais, a máscara de alvo observado permite ignorar pontos imputados da função de perda e das métricas. Assim, a avaliação de teste é comparável, mas o mecanismo de treino não é perfeitamente simétrico.

Além da comparação principal, foram executados experimentos diagnósticos com o Seq2Seq atenção para investigar limites informacionais do problema. O primeiro é o \textit{oracle future auxiliary}: no decodificador, o modelo recebe PM10 e partículas totais em suspensão (PTS) reais do próprio horizonte futuro. Esse cenário não representa um uso causal de previsão, pois usa informação indisponível no instante de previsão, mas mede quanto desempenho poderia ser recuperado caso variáveis auxiliares futuras relevantes fossem conhecidas. O segundo é o Seq2Seq multi-alvo, no qual a rede continua prevendo PM2,5 como alvo principal, mas recebe uma cabeça auxiliar para prever PM10 e PTS futuros durante o treinamento. Esse teste avalia se aprender variáveis correlacionadas em conjunto ajuda a representação interna sem introduzir vazamento na avaliação causal.

Esses dois experimentos não substituem o protocolo final nem entram como candidatos a vencedor, porque usam contratos diferentes do contrato causal principal. Eles são apresentados nos resultados como análise de limite: o \textit{oracle} mostra o ganho máximo associado a variáveis futuras observadas, enquanto o multi-alvo testa se a própria arquitetura consegue aproximar esse ganho apenas aprendendo as auxiliares.

\section{Treinamento e Busca de Hiperparâmetros}

O resultado principal usa hiperparâmetros selecionados no bloco de desenvolvimento anterior ao teste. Para LSTM direta, LSTM recursiva e XGBoost, a seleção foi feita por validação temporal \textit{walk-forward}. Para os modelos neurais principais com HPO, foram usadas 24 tentativas, três partições temporais por janela expansiva, até 18 épocas por tentativa e treinamento final de até 36 épocas para LSTM direta e LSTM recursiva. Para o novo Seq2Seq atenção, a configuração final veio de ablações focadas em arquitetura, perda, resposta a picos e tamanho de janela, e foi reexecutada no mesmo protocolo final com quatro sementes. O XGBoost foi submetido ao mesmo desenho \textit{walk-forward}, mas, por custo computacional das tentativas mais longas, o resultado consolidado usou quatro tentativas completas.

Nas redes neurais, cada tentativa reiniciou o modelo em cada partição e agregou a métrica de validação ao fim das partições. O estudo do Optuna usou o amostrador padrão \texttt{TPESampler} e o \texttt{HyperbandPruner}, com \texttt{min\_resource}=1, \texttt{max\_resource}=18 e fator de redução 3. Assim, tentativas pouco promissoras podiam ser interrompidas antes de consumir todo o orçamento de épocas. Essa escolha segue a recomendação da documentação do Optuna para a combinação entre \texttt{TPESampler} e \texttt{HyperbandPruner} \cite{optunaDocsEfficient2026,optunaDocsHyperband2026}. Para o XGBoost, a busca minimizou o MAE médio nas partições, com parada antecipada (\textit{early stopping}) nativa do XGBoost e \texttt{MedianPruner} \cite{optunaDocsMedian2026}. A Tabela \ref{tab:hpo_espaco_busca} organiza o espaço de busca por grupo de hiperparâmetros; a Tabela \ref{tab:hpo_melhores_hiperparametros} apresenta a configuração selecionada.

Nos modelos neurais, a taxa de aprendizado apresentada nas tabelas é a taxa nominal \(\eta_{\max}\) escolhida pela HPO. As execuções finais usaram AdamW com agenda ativa por época: três épocas de aquecimento linear, saindo de \(10^{-5}\eta_{\max}\) até \(\eta_{\max}\), seguidas de decaimento cosseno até \(0{,}01\eta_{\max}\) no fim do orçamento. A agenda é reiniciada em cada partição da HPO e em cada treinamento final. A validação ao fim das épocas monitora MAE para poda, parada antecipada e escolha do melhor \textit{checkpoint}, mas não recalcula a taxa de aprendizado com base no valor da métrica.

\begin{table}[htb]
    \centering
    \caption{Espaço de busca e configuração selecionada na comparação principal.}
    \label{tab:hpo_espaco_busca}
    \small
    \renewcommand{\arraystretch}{1.08}
    \begin{tabular}{>{\raggedright\arraybackslash}p{0.19\textwidth}>{\raggedright\arraybackslash}p{0.23\textwidth}>{\raggedright\arraybackslash}p{0.48\textwidth}}
        \toprule
        \textbf{Grupo} & \textbf{Hiperparâmetro} & \textbf{Valores avaliados} \\
        \midrule
        Redes neurais & Taxa nominal de aprendizado & 0{,}0003; 0{,}0005; 0{,}0008; 0{,}0012; 0{,}0018; 0{,}0025; 0{,}0035; 0{,}005; 0{,}007 \\
        Redes neurais & Decaimento de pesos & 0; $10^{-7}$; $5 \times 10^{-7}$; $10^{-6}$; $5 \times 10^{-6}$; $10^{-5}$; $5 \times 10^{-5}$; $10^{-4}$ \\
        Redes neurais & Taxa de \textit{dropout} & 0{,}00; 0{,}05; 0{,}10; 0{,}15; 0{,}20; 0{,}25; 0{,}30 \\
        Redes neurais & Profundidade e perda & número de camadas 1; 2; 3; $\delta$ da Huber em 0{,}5; 0{,}75; 1{,}0; 1{,}5; 2{,}0 \\
        LSTM direta & Arquitetura & dimensão oculta 96; 128; 192; 256; 320; agregação por último estado, média ou atenção \\
        LSTM recursiva & Arquitetura & dimensão oculta 96; 128; 192; 256; 320 \\
        Seq2Seq atenção novo & Arquitetura e perda & variante selecionada por ablações: codificador/decodificador 96; uma camada; atenção de entrada 96; representação de horizonte 24; reinício do decodificador por horizonte; sem \textit{input feeding}; atenção por \textit{lag buckets}; perda ponderada e perda auxiliar leve de evento \\
        XGBoost com saída múltipla & Árvores e profundidade & 160; 240; 320; 480 árvores; profundidade máxima 2 a 7; \texttt{max\_bin} 64; 128; 256 \\
        XGBoost com saída múltipla & Regularização, amostragem e parada & taxa de aprendizado logarítmica entre 0{,}015 e 0{,}12; \textit{subsample} e \textit{colsample} entre 0{,}60 e 1{,}0; $\lambda$, $\alpha$, $\gamma$, peso mínimo do filho (\textit{min child weight}); parada antecipada (\textit{early stopping}) de 10, 15 ou 25 rodadas \\
        \bottomrule
    \end{tabular}
\end{table}

\begin{table}[htb]
    \centering
    \caption{Hiperparâmetros e configurações selecionadas na comparação principal.}
    \label{tab:hpo_melhores_hiperparametros}
    \small
    \renewcommand{\arraystretch}{1.08}
    \begin{tabular}{>{\raggedright\arraybackslash}p{0.20\textwidth}>{\raggedright\arraybackslash}p{0.35\textwidth}>{\raggedright\arraybackslash}p{0.35\textwidth}}
        \toprule
        \textbf{Modelo} & \textbf{Arquitetura selecionada} & \textbf{Treinamento e regularização} \\
        \midrule
        LSTM direta & dimensão oculta 256; 3 camadas; agregação por média & taxa nominal de aprendizado 0{,}007; decaimento de pesos $5 \times 10^{-5}$; \textit{dropout} 0{,}20; $\delta$ da Huber 2{,}0 \\
        LSTM recursiva & dimensão oculta 128; 2 camadas & taxa nominal de aprendizado 0{,}0035; decaimento de pesos 0; \textit{dropout} 0{,}20; $\delta$ da Huber 1{,}0 \\
        Seq2Seq atenção novo & codificador 96; decodificador 96; 1 camada; atenção de entrada 96; representação de horizonte 24; \textit{lag buckets} com janela 2 & taxa nominal de aprendizado 0{,}005; decaimento de pesos $10^{-6}$; \textit{dropout} 0{,}25; pesos 1{,}05/1{,}10; fator de subestimação 1{,}50; perda auxiliar de evento 0{,}035 \\
        XGBoost com saída múltipla & 480 árvores; profundidade máxima 2; \texttt{max\_bin}=64 & taxa de aprendizado 0{,}02563; \textit{subsample}=0{,}7611; \textit{colsample}=0{,}9073; $\lambda=7{,}7542$; parada antecipada=10 rodadas \\
        \bottomrule
    \end{tabular}
\end{table}

\clearpage

A formulação final do Seq2Seq priorizou a nova variante com atenção de entrada, perda \texttt{weighted\_l1}, reinício do estado do decodificador a cada horizonte, remoção de \textit{input feeding}, atenção condicionada por horizonte e variáveis do decodificador, agrupamento por \textit{lag buckets} e perda auxiliar leve para eventos. A perda pondera regimes de concentração e tipos de erro, alterando o objetivo de treinamento para reduzir parte da suavização e melhorar o comportamento em concentrações mais altas. Essa variante foi avaliada no protocolo consolidado com múltiplas sementes e janela fixa de 48 horas, após ablações indicarem que janelas de 72, 96 e 168 horas não traziam ganho claro.

As variantes finais de Seq2Seq priorizaram um decodificador semidireto com horizonte explícito e linha de base residual sazonal, reduzindo a dependência de realimentação autoregressiva pura. Ainda assim, a avaliação de agendas de \textit{teacher forcing} foi incluída porque trabalhos recentes em previsão, como \citeonline{teutsch2022flipped}, indicam que currículos crescentes de \textit{teacher forcing} podem melhorar previsões multi-horizonte em séries temporais contínuas. No relatório, portanto, \textit{teacher forcing} e \textit{scheduled sampling} são tratados como técnicas avaliadas em ablações auxiliares fundamentadas, não como eixo central do resultado final em Sapo. Essa distinção evita atribuir ao \textit{scheduled sampling} um ganho que, nos experimentos consolidados, esteve mais associado ao contrato de dados, ao uso causal de PM10, à seleção por validação e à função de perda.

\section{Métricas de Avaliação}

A avaliação final utiliza métricas no bloco final de teste. As métricas globais principais são erro médio absoluto (MAE), raiz do erro quadrático médio (RMSE), erro percentual médio absoluto (MAPE) e coeficiente de determinação (\(R^2\)). Também são reportadas métricas por horizonte, especialmente nos horizontes 1, 6, 12 e 24, para observar degradação ao longo da previsão; os rótulos H1, H6, H12 e H24 indicam esses horizontes em horas. Como os picos são raros e importantes em qualidade do ar, a análise também considera métricas de cauda, como erro em eventos acima de 35, percentil 99 previsto (p99), razão entre desvio-padrão previsto e observado, viés médio medido pelo erro médio de viés (\textit{Mean Bias Error}, MBE) e taxa de subestimação. Nas tabelas, o rótulo Peak35 MAE indica o MAE calculado em pontos ou eventos com PM2,5 \(\geq 35\).

A interpretação dos resultados combina métricas agregadas e inspeção visual das séries previstas. Essa combinação é necessária porque modelos que vencem em erro médio podem produzir previsões mais suaves, enquanto variantes voltadas a picos podem preservar mais amplitude ao custo de pequeno aumento no MAE global.

\section{Reprodutibilidade}

A reprodutibilidade foi tratada por meio de scripts determinísticos de geração de figuras, tabelas e diagnósticos, mantendo a separação entre treino, validação e teste. Os experimentos principais foram executados com identificadores de execução, sementes registradas quando aplicável, configuração explícita de HPO e exportação de métricas em arquivos estruturados. No corpo do TCC, os valores reportados são apresentados diretamente nas tabelas e nas figuras; os caminhos técnicos e hashes dos arquivos ficam preservados como material de controle do projeto, sem substituir a descrição metodológica.
